% Provenance (never printed): see sources/outline_sources.md and sources/book1/BRIEFS.md (g6/05).
\chapter{Repeating with \texttt{while}}\label{ch:g6:while-loops}

Tom wants a game that costs 50 coins. He saves 3 coins every week. How many
weeks must he wait? He could work it out, but he could also just count: one
week, 3 coins; two weeks, 6 coins; and so on, \emph{as long as} he has fewer
than 50 coins. A \code{for} \omterm{def:g4:algorithms-repeat-and-decide:loop}{loop} needs to know in advance how many times to
turn, and here that number is exactly what we are looking for.

\begin{recall}
A \code{for} \omterm{def:g4:algorithms-repeat-and-decide:loop}{loop} repeats its block a known number of times
(\cref{ch:g6:for-loops-turtle}). A \omterm{def:g6:if-else:comparison}{comparison} such as \code{coins < 50} gives
\code{True} or \code{False} (\cref{ch:g6:if-else}). An \omterm{def:g4:algorithms-repeat-and-decide:algorithm}{algorithm} can repeat
steps until a \omterm{def:g4:algorithms-repeat-and-decide:condition}{condition} becomes true (\cref{ch:g4:algorithms-repeat-and-decide}).
\end{recall}

\section{Repeating as long as a condition holds}

\begin{definition}[While loop]\label{def:g6:while-loops:while}
A \emph{while loop}\index{while loop} repeats its block as long as a \omterm{def:g4:algorithms-repeat-and-decide:condition}{condition}
is true: \code{while condition:} followed by an indented block. The \omterm{def:g4:algorithms-repeat-and-decide:condition}{condition}
is checked before every turn; as soon as it is false, the \omterm{def:g4:algorithms-repeat-and-decide:loop}{loop} stops and the
\omterm{def:g4:programs-and-operating-system:program}{program} goes on after the block.
\end{definition}

\Needspace*{8\baselineskip}
\omcode{code/grade-6/05-while-loops/listings/countdown.py}{1}{5}{A
countdown.}\label{lst:g6:while-loops:countdown}

\Needspace*{10\baselineskip}
\omoutput{code/grade-6/05-while-loops/outputs/countdown.txt}{When \code{n}
reaches 0, the \omterm{def:g4:algorithms-repeat-and-decide:condition}{condition} \code{n > 0} is false and the \omterm{def:g4:algorithms-repeat-and-decide:loop}{loop}
stops.}\label{out:g6:while-loops:countdown}

\begin{omfigure}
\begin{tikzpicture}[x=1cm, y=1cm, font=\footnotesize,
    q/.style={draw=black!70, diamond, aspect=2.2, fill=ocYellow!35, inner sep=1pt},
    blk/.style={draw=black!65, rounded corners=3pt, fill=#1, minimum width=2.6cm, minimum height=0.75cm, align=center}]
  \node[blk=black!6] (s) at (0,2.2) {\code{n = 5}};
  \node[q] (c) at (0,0.8) {\code{n > 0}?};
  \node[blk=ocBlue!12] (b) at (0,-0.9) {\code{print(n)}\\\code{n = n - 1}};
  \node[blk=black!6] (e) at (4.3,0.8) {\code{print("Lift-off!")}};
  \draw[-{Stealth[length=2mm]}, thick] (s) -- (c);
  \draw[-{Stealth[length=2mm]}, thick] (c) -- node[right] {\code{True}} (b);
  \draw[-{Stealth[length=2mm]}, thick] (b.west) -- ++(-0.7,0) |- (c.west);
  \draw[-{Stealth[length=2mm]}, thick] (c) -- node[above] {\code{False}} (e);
  \node[black!60, align=right, anchor=east] at (-2.15,-0.05) {back to\\the test};
\end{tikzpicture}

{\small A \code{while} \omterm{def:g4:algorithms-repeat-and-decide:loop}{loop} is a cycle with the test at its door: the block
runs, then the \omterm{def:g4:programs-and-operating-system:program}{program} goes back to the test. If the test is false the very
first time, the block never runs at all.}
\end{omfigure}

\section{Counters and running totals}

Tom's savings need two \omterm{def:g6:variables-and-input:value}{variables} that change at each turn: the number of coins
and the number of weeks.

\begin{definition}[Counter and running total]\label{def:g6:while-loops:counter}
A \emph{counter}\index{counter} is a \omterm{def:g6:variables-and-input:value}{variable} that starts at 0 and goes up by
1 each time something happens. A \emph{running total}\index{running total} is
a \omterm{def:g6:variables-and-input:value}{variable} that starts at 0 and adds each new \omterm{def:g6:variables-and-input:value}{value} as it comes.
\end{definition}

\Needspace*{9\baselineskip}
\omcode{code/grade-6/05-while-loops/listings/savings.py}{1}{6}{Tom's
savings.}\label{lst:g6:while-loops:savings}

\Needspace*{5\baselineskip}
\omoutput{code/grade-6/05-while-loops/outputs/savings.txt}{After 16 weeks Tom has
48 coins, still fewer than 50; after 17 weeks he has
51.}\label{out:g6:while-loops:savings}

\begin{omfigure}
\begin{tikzpicture}[x=0.16cm, y=1cm, font=\scriptsize]
  \foreach \w in {1,...,17} {
    \pgfmathtruncatemacro{\c}{3*\w}
    \fill[ocYellow!70!black!60] (0,-0.32*\w) rectangle (\c,-0.32*\w+0.24);
    \node[anchor=east] at (-0.6,-0.32*\w+0.12) {week \w};
    \node[anchor=west] at (\c+0.4,-0.32*\w+0.12) {\c};
  }
  \draw[ocRed, thick, dashed] (50,0.1) -- (50,-5.6) node[below] {50 coins};
\end{tikzpicture}

{\small The trace of \cref{lst:g6:while-loops:savings} as a filling bar: one row
per turn of the \omterm{def:g4:algorithms-repeat-and-decide:loop}{loop}. The test \code{coins < 50} is still true after week 16 (48
coins), and false after week 17 (51 coins).}
\end{omfigure}

A \omterm{def:g6:while-loops:counter}{running total} can also add numbers typed by the user. Since we do not know
how many numbers will come, we agree on a special \omterm{def:g6:variables-and-input:value}{value} that means ``stop'':
here, 0. The \omterm{def:g4:programs-and-operating-system:program}{program} reads one number before the \omterm{def:g4:algorithms-repeat-and-decide:loop}{loop}, and the next one at the
end of each turn, so the \omterm{def:g4:algorithms-repeat-and-decide:condition}{condition} always tests the latest number.

\Needspace*{9\baselineskip}
\omcode{code/grade-6/05-while-loops/listings/sum_until_zero.py}{1}{6}{Adding
typed numbers until 0.}\label{lst:g6:while-loops:sum}

\Needspace*{9\baselineskip}
\omoutput{code/grade-6/05-while-loops/outputs/sum_until_zero.txt}{$12 + 7 + 30
= 49$; the 0 only stops the \omterm{def:g4:algorithms-repeat-and-decide:loop}{loop}.}\label{out:g6:while-loops:sum}

With a \omterm{def:g6:while-loops:counter}{counter} as well, the \omterm{def:g4:programs-and-operating-system:program}{program} can compute an average: the total divided
by the number of \omterm{def:g6:variables-and-input:value}{values}. A mark can be 0, so here the stop \omterm{def:g6:variables-and-input:value}{value} is $-1$.

\Needspace*{11\baselineskip}
\omcode{code/grade-6/05-while-loops/listings/average.py}{1}{8}{The average of
typed marks.}\label{lst:g6:while-loops:average}

\Needspace*{10\baselineskip}
\omoutput{code/grade-6/05-while-loops/outputs/average.txt}{$(14 + 9.5 + 17 +
12.5) \div 4 = 13.25$.}\label{out:g6:while-loops:average}

\section{The loop that never ends}

\begin{definition}[Infinite loop]\label{def:g6:while-loops:infinite}
An \emph{infinite loop}\index{infinite loop} is a \omterm{def:g4:algorithms-repeat-and-decide:loop}{loop} whose \omterm{def:g4:algorithms-repeat-and-decide:condition}{condition} never
becomes false: the \omterm{def:g4:programs-and-operating-system:program}{program} keeps turning and never reaches the lines after the
\omterm{def:g4:algorithms-repeat-and-decide:loop}{loop}.
\end{definition}

Tom first wrote his savings \omterm{def:g4:programs-and-operating-system:program}{program} with \code{!=} instead of \code{<}: ``as
long as I do not have exactly 50 coins''. But 3 coins a week give 48, then 51:
never exactly 50.

\Needspace*{7\baselineskip}
\omcode{code/grade-6/05-while-loops/listings/forever.py}{2}{5}{Waiting for
exactly 50 coins.}\label{lst:g6:while-loops:forever}

\Needspace*{6\baselineskip}
\omoutput{code/grade-6/05-while-loops/outputs/forever.txt}{The \omterm{def:g4:programs-and-operating-system:program}{program} prints
nothing and never ends. In this book, a \omterm{def:g4:programs-and-operating-system:program}{program} that runs too long is stopped
after 2 seconds.}\label{out:g6:while-loops:forever}

On a \omterm{def:g1:what-is-a-computer:computer}{computer}, a \omterm{def:g4:programs-and-operating-system:program}{program} caught in an \omterm{def:g6:while-loops:infinite}{infinite loop} is stopped by pressing
\key{Ctrl} and \key{C} together in its terminal. To fix the \omterm{def:g4:algorithms-repeat-and-decide:loop}{loop}, check two
things: the block must change a \omterm{def:g6:variables-and-input:value}{variable} of the \omterm{def:g4:algorithms-repeat-and-decide:condition}{condition}, and the change must
lead the \omterm{def:g4:algorithms-repeat-and-decide:condition}{condition} to become false. Here, \code{coins < 50} becomes false as
soon as the 50 coins are reached or passed:

\Needspace*{7\baselineskip}
\omcode{code/grade-6/05-while-loops/listings/fixed_forever.py}{1}{4}{The
\omterm{def:g4:algorithms-repeat-and-decide:condition}{condition} becomes false.}\label{lst:g6:while-loops:fixed}

\Needspace*{5\baselineskip}
\omoutput{code/grade-6/05-while-loops/outputs/fixed_forever.txt}{The \omterm{def:g4:algorithms-repeat-and-decide:loop}{loop}
stops.}\label{out:g6:while-loops:fixed}

\section{Choosing between \texorpdfstring{\texttt{for}}{for} and \texorpdfstring{\texttt{while}}{while}}

Any \code{for} \omterm{def:g4:algorithms-repeat-and-decide:loop}{loop} over a \code{range} can be written with \code{while} and a
\omterm{def:g6:while-loops:counter}{counter}, but it is longer, and forgetting the line that changes the \omterm{def:g6:while-loops:counter}{counter}
gives an \omterm{def:g6:while-loops:infinite}{infinite loop}.

\Needspace*{10\baselineskip}
\omcode{code/grade-6/05-while-loops/listings/for_and_while.py}{1}{7}{The same
countdown, both ways. The \code{while} version must start, test and change
\code{n} itself.}\label{lst:g6:while-loops:for-and-while}

\Needspace*{10\baselineskip}
\omoutput{code/grade-6/05-while-loops/outputs/for_and_while.txt}{The two \omterm{def:g4:algorithms-repeat-and-decide:loop}{loops} print
the same numbers.}\label{out:g6:while-loops:for-and-while}

\begin{method}[Choosing a loop]\label{met:g6:while-loops:choose}
If the number of turns is known before the \omterm{def:g4:algorithms-repeat-and-decide:loop}{loop} starts (``for each of the 5
squares'', ``for every age from 0 to 80''), use \code{for}. If the \omterm{def:g4:algorithms-repeat-and-decide:loop}{loop} must go
on until something happens (``until 0 is typed'', ``as long as the guess is
wrong''), use \code{while}, and check that its \omterm{def:g4:algorithms-repeat-and-decide:condition}{condition} will become false.
\end{method}

How many digits does a whole number have? We do not know in advance, but we can
remove its last digit with \code{// 10} until a single digit is left, and count:

\Needspace*{9\baselineskip}
\omcode{code/grade-6/05-while-loops/listings/digits.py}{1}{6}{Counting
digits.}\label{lst:g6:while-loops:digits}

\Needspace*{6\baselineskip}
\omoutput{code/grade-6/05-while-loops/outputs/digits.txt}{2026, 202, 20, 2: three
turns, plus the last digit.}\label{out:g6:while-loops:digits}

\section{A number-guessing game}

Two players: the first types a secret number from 1 to 100, which the \omterm{def:g4:programs-and-operating-system:program}{program}
then pushes out of sight with three empty lines; the second guesses, and the
\omterm{def:g4:programs-and-operating-system:program}{program} answers ``higher'' or ``lower'' until the guess is right, counting the
guesses.

\omcode{code/grade-6/05-while-loops/listings/guess.py}{1}{12}{A two-player
guessing game.}\label{lst:g6:while-loops:guess}

\omoutput{code/grade-6/05-while-loops/outputs/guess.txt}{The secret was
37.}\label{out:g6:while-loops:guess}

A good player always guesses in the middle of the numbers still possible: each
answer then removes half of them.

\begin{omfigure}
\begin{tikzpicture}[x=0.11cm, y=1cm, font=\scriptsize]
  \foreach \lo/\hi/\g/\r/\t [count=\i] in {1/100/50/0/{guess 50: lower}, 1/49/25/0/{guess 25: higher}, 26/49/37/1/{guess 37: found}} {
    \draw[black!40] (1,-0.8*\i) rectangle (100,-0.8*\i+0.35);
    \fill[ocBlue!35] (\lo,-0.8*\i) rectangle (\hi,-0.8*\i+0.35);
    \ifnum\r=1 \fill[ocMint!80!black] (\g-0.5,-0.8*\i) rectangle (\g+0.5,-0.8*\i+0.35);
    \else \draw[ocRed, thick] (\g,-0.8*\i-0.05) -- (\g,-0.8*\i+0.4); \fi
    \node[anchor=west] at (103,-0.8*\i+0.17) {\lo\ to \hi: \t};
  }
  \foreach \x in {1,50,100} \node[below] at (\x,-2.45) {\x};
\end{tikzpicture}

{\small The numbers still possible (blue) before each guess (red line), in the
game above. Halving the possible numbers, 7 guesses are always enough for 1 to
100.}
\end{omfigure}

\section{Exercises}

\begin{exercise}[$\star$]\label{exo:g6:while-loops:1}
What does this \omterm{def:g4:programs-and-operating-system:program}{program} print?

\omcode{code/grade-6/05-while-loops/exercises/by_two.py}{1}{4}{Counting down by
two.}\label{lst:g6:while-loops:by-two}
\end{exercise}

\begin{exercise}[$\star$]\label{exo:g6:while-loops:2}
How many times does the block of this \omterm{def:g4:algorithms-repeat-and-decide:loop}{loop} run? What is the last number
printed?

\omcode{code/grade-6/05-while-loops/exercises/turns.py}{1}{4}{Count the
turns.}\label{lst:g6:while-loops:turns}
\end{exercise}

\begin{exercise}[$\star$]\label{exo:g6:while-loops:3}
What does this \omterm{def:g4:programs-and-operating-system:program}{program} print? Make its \omterm{def:g6:variables-and-input:trace-table}{trace table}.

\omcode{code/grade-6/05-while-loops/exercises/squares_total.py}{1}{6}{A \omterm{def:g6:while-loops:counter}{running
total}.}\label{lst:g6:while-loops:squares-total}
\end{exercise}

\begin{exercise}[$\star$]\label{exo:g6:while-loops:4}
Write a \omterm{def:g4:programs-and-operating-system:program}{program} that adds the whole numbers from 1 to 100 with a \code{while}
\omterm{def:g4:algorithms-repeat-and-decide:loop}{loop}.
\end{exercise}

\begin{exercise}[$\star\star$]\label{exo:g6:while-loops:5}
This \omterm{def:g4:programs-and-operating-system:program}{program} never stops. Why? Fix it so that it prints the numbers from 0 to 9.

\omcode{code/grade-6/05-while-loops/exercises/never_bug.py}{1}{3}{A \omterm{def:g4:algorithms-repeat-and-decide:loop}{loop} that
never ends.}\label{lst:g6:while-loops:never-bug}
\end{exercise}

\begin{exercise}[$\star\star$]\label{exo:g6:while-loops:6}
This \omterm{def:g4:programs-and-operating-system:program}{program} should print $1 + 2 + \dots + 10 = 55$, but prints 45. Which \omterm{def:g6:variables-and-input:value}{value}
is missing, and why? Fix it.

\omcode{code/grade-6/05-while-loops/exercises/total_bug.py}{1}{7}{A total that
misses a \omterm{def:g6:variables-and-input:value}{value}.}\label{lst:g6:while-loops:total-bug}
\end{exercise}

\begin{exercise}[$\star\star$]\label{exo:g6:while-loops:7}
Starting from 1 and doubling again and again, which is the first number above
1,000? Write a \omterm{def:g4:programs-and-operating-system:program}{program} that finds it and counts the doublings.
\end{exercise}

\begin{exercise}[$\star\star$]\label{exo:g6:while-loops:8}
Write a \omterm{def:g4:programs-and-operating-system:program}{program} that asks for a whole number and prints the sum of its digits:
2026 gives $2 + 0 + 2 + 6 = 10$. Use \texttt{\% 10} to get the last digit.
\end{exercise}

\begin{exercise}[$\star\star$]\label{exo:g6:while-loops:9}
How many turns does \code{while n > 1: n = n // 2} take when \code{n} starts at
100? Make the \omterm{def:g6:variables-and-input:trace-table}{trace table}.
\end{exercise}

\begin{exercise}[$\star\star\star$]\label{exo:g6:while-loops:10}
Write a \omterm{def:g4:programs-and-operating-system:program}{program} that asks ``Have you finished your homework?'' again and again,
until the answer is \code{yes}.
\end{exercise}

\begin{exercise}[$\star\star\star$]\label{exo:g6:while-loops:11}
In \cref{lst:g6:while-loops:savings}, replace the 3 coins a week by any whole
number of coins greater than 0. Explain why the \omterm{def:g4:algorithms-repeat-and-decide:loop}{loop} always stops. Would it
still stop with 0 coins a week?
\end{exercise}

\begin{exercise}[$\star\star\star$]\label{exo:g6:while-loops:12}
Change the guessing game so that player 2 has at most 7 guesses; after the
seventh wrong guess, the \omterm{def:g4:programs-and-operating-system:program}{program} says that player 1 wins and shows the secret.
\end{exercise}

\begin{history}[An unsolved problem]
% ledger: g6:collatz-author, g6:collatz-open
The game of the weekend problem below is often called the Collatz problem,
after the mathematician Lothar Collatz, who wrote down questions of this
kind in his notebooks in the 1930s.
It remains unsolved: nobody has proved that every start comes down to 1, and
nobody has found a start that does not.
\end{history}

\section{Problem: The 3n + 1 Game}

\begin{problem}[{Weekend problem --- does every number come down to 1?}]\label{pb:g6:while-loops:1}
Start from a whole number. If it is even, halve it; if it is odd, triple it and
add 1. Do the same with the result, again and again, and stop when you reach 1.
Each change is one \emph{step}. From 3, for instance: 3, 10, 5, 16, 8, 4, 2, 1,
seven steps.

\medskip\noindent\textbf{Part I --- By hand.}
\begin{enumerate}
  \item Play the game from 6. How many steps does it take? What is the
    highest number reached?
  \item Play it from 7. How many steps, and what is the highest number?
  \item Why does the game end at once when the number is a doubling of 1, such
    as 16, 32 or 64? How many steps does 64 take?
\end{enumerate}

\medskip\noindent\textbf{Part II --- A program.}

\omcode{code/grade-6/05-while-loops/listings/collatz.py}{1}{12}{Steps and highest
number for a typed start.}\label{lst:g6:while-loops:collatz}

\begin{enumerate}[resume]
  \item Which \omterm{def:g6:variables-and-input:value}{variables} of \cref{lst:g6:while-loops:collatz} are \omterm{def:g6:while-loops:counter}{counters}, and
    which keeps a record? Why must \code{highest} start at the starting number?
  \item Run the \omterm{def:g4:programs-and-operating-system:program}{program} for 6 and 7. Does it agree with your answers?
  \item Run it for 27. What do you notice?
  \item Write a \omterm{def:g4:programs-and-operating-system:program}{program} that plays the game from every start from 1 to 99 and
    prints the start that takes the most steps, with its number of steps. (Put
    the \code{while} \omterm{def:g4:algorithms-repeat-and-decide:loop}{loop} inside a \code{for} \omterm{def:g4:algorithms-repeat-and-decide:loop}{loop}.)
\end{enumerate}

\medskip\noindent\textbf{Part III --- An open question.}
\begin{enumerate}[resume]
  \item For the start you found in question 7, what is the highest number
    reached?
  \item Your \omterm{def:g4:programs-and-operating-system:program}{program} of question 7 stopped for every start below 100. Does
    that prove that every start reaches 1?
  \item What would happen to your \omterm{def:g4:programs-and-operating-system:program}{program} if some start never reached 1?
  \item Mathematicians have played this game with \omterm{def:g1:what-is-a-computer:computer}{computers} for every start
    up to enormous numbers, and all reached 1. Yet nobody knows whether every
    start does. Explain the difference between ``checked for many numbers'' and
    ``proved for all numbers''.
  \item State the result of question 7.
\end{enumerate}
\end{problem}
